
% both accelerate and broader our understanding of future infectious disease outbreaks

% Pandemic preparedness is a cornerstone of public health safety. In addition logistic and governance challenges to facilitate research dissemination and coordinate international efforts, its scientific underpinnings are incredibly multi-faceted. As illustrated by the recent COVID-19 pandemic, we don't only need large-scale genomic surveillance to enhance phylodynamic analyses (good data), but many other data sources to get a more complete picture of transmission dynamics, the biological mechanisms of pathogen-receptor binding, but also early detection of infection in potential reservoir. We also need good methods i.e. fast computing methods and efficient data engineering. In that light, the primary objective of this thesis is to explore how new efficient, vector-based, representations of various datasets --- phylogenetic trees, protein structures, and animal tracking --- can better our understanding of pandemics. First, I developed phylo2vec, an integer vector-based representation of binary phylogenetic trees. Designed with rapid sampling and efficient data storage in mind, this compact representation also holds potential for discrete, hierarchical phylogenetic data amenable for use in downstream machine learning analyses. Subsequently, I leveraged autoencoder-based vector representations of protein structure to assess how recombination and mutation shape receptor affinity in \textit{Betacoronavirus} spike glycoproteins, thus complementing traditional sequence-based phylogenetic analysis with structural insights. This analysis refines previous models of \textit{Betacoronavirus} protein evolution, suggesting that while recombination events and gradual mutations in the N-terminal domain (NTD) appear to significantly drive ACE2 binding capability at the sequence level, receptor switching is mediated by gradual structural refinement within the receptor-binding domain (RBD) rather than by large-scale structural reorganisation. Lastly, I developed a feasibility study to assess the potential of deep learning-based tracking for the early detection of infection in broiler chickens. Using \textit{E. coli}, a major pathogen of concern in poultry farming, we were able to identify behaviour patterns associated with impaired welfare from videography, suggesting a promising avenue for the use of artificial intelligence, by turning video surveillance data into a vector of insightful behavioural features, for modelling and monitoring infectious disease epidemics in poultry. Together, this thesis illustrates the inhernet cross-disicplinary of pandemic preparedness. Specifically, ti demonstrates how better data engineering, as well as more comprehensive data collection beyond the now traditional genomic surveillance, may both accelerate and broader our understanding of future infectious disease outbreaks.


% influence of recombination and mutation in mediating the affinity of spike glycoproteins from betacoronaviruses with angiotensin-converting enzyme 2.

% new data representations for phylogenetic trees to make them amenable 


% multi-faceted and difficult. Beyond logistic and governance challenges and scientific challenges due to advent of big genomic data since the COVID-19 pandemic which has shown that collecting data is beneficial for analyses. 